\chapter{Corrigés des exercices}\label{annexe:corriges}
\iterbanner{V3}

\section*{Chapitre~\ref{ch:constat} — Le constat}
\textbf{Ex.~\ref{ex:constat-budget}.} (a) $300\times0{,}43=129$ profils faisables. (b) Les pertes passent de 57~\% à 28,5~\% : $300\times0{,}715\approx215$ faisables, ce que l'approche actuelle n'obtiendrait qu'avec $215/0{,}43\approx500$ simulations. (c) Le front est au bord des contraintes : une méthode qui le cherche propose forcément des profils de part et d'autre de la frontière ; les modèles appris sur quelques centaines de points se trompent près de la frontière ; et la convergence d'XFOIL n'est connue qu'après simulation. Le hasard fait 76~\% parce qu'il ne cherche pas le bord.

\section*{Chapitre~\ref{ch:taxonomie} — Classer les contraintes}
\textbf{Ex.~\ref{ex:taxo-classer}.} (a) Envergure : connue, a priori, quantifiable, relâchable (le code calcule au-delà) ; on l'impose par les bornes de la grille, sans pénalité. (b) Mailleur qui échoue : cachée, non quantifiable, non relâchable ; mais le mailleur est rapide, donc on le lance \emph{avant} le solveur, comme un filtre, et l'on apprend un classifieur pour orienter. (c) Flutter : connue, issue de la simulation, quantifiable (la marge), relâchable ; un modèle de substitution à part, comme $C_{L,\max}$. (d) Solution instationnaire : connue mais non quantifiable si l'on n'a qu'un indicateur, quantifiable si l'on mesure l'amplitude des oscillations ; issue de la simulation ; les sorties existent mais n'ont pas de sens (non relâchable en pratique) : un classifieur, ou un modèle de l'amplitude.

\textbf{Ex.~\ref{ex:taxo-filtre}.} Budget : une pénalité finie, sur un modèle approché, n'empêche pas de proposer un profil qui la viole ; chaque proposition de ce type est une simulation perdue, et comme la contrainte n'est pas relâchable, elle ne rend même pas d'information. Modèle : les profils filtrés n'entrent jamais dans les données ; avec une pénalité, la boucle simule ces profils, ou au contraire n'apprend rien près de la frontière, et le modèle extrapole. Le filtre est exact et gratuit, la pénalité est approchée et coûteuse.

\section*{Chapitre~\ref{ch:penalites} — Pénalités}
\textbf{Ex.~\ref{ex:pen-rentable}.} (a) Le profil faisable vaut $-0{,}80$. Le voisin vaut $-1{,}00+75\times0{,}04^2=-1{,}00+0{,}12=-0{,}88$ : c'est lui, infaisable, l'optimum du QUBO. (b) Il faut $P\times0{,}0016>0{,}20$, soit $P>125$. (c) Pour $h=-0{,}5$, la pénalité vaut $125\times0{,}25\approx31$, pour un objectif d'étendue 1 : l'objectif n'occupe plus que 3~\% environ du spectre. Rendre la pénalité sûre pour les petites violations la rend écrasante pour les grandes.

\textbf{Ex.~\ref{ex:pen-deseq}.} (a) $75(h^2-0{,}1h)$ : $0{,}56$ en $h=-0{,}05$ ; $0$ en $h=0$ ; $-0{,}19$ en $h=0{,}05$ ; $4{,}5$ en $h=0{,}3$. (b) Pour $h>2h^\star=0{,}1$ : tout profil dont la marge dépasse 0,1 est pénalisé. (c) Avec $h^\star=0{,}15$, la pénalité s'annule en $h=0{,}3$ quel que soit $P$ ; avec $P=30$, elle vaut $0{,}53$ en $h=-0{,}05$ et $-0{,}68$ en $h=0{,}15$. On ne pénalise plus les profils de marge $0{,}3$, mais on attire maintenant les candidats vers une marge de 0,15, loin du front qui longe la contrainte. On ne peut pas avoir les deux avec une parabole.

\textbf{Ex.~\ref{ex:pen-degre}.} Avec $z_i=2b_i-1$ : $h=4b_1b_2-2b_1-2b_2+0{,}8$. Dans $h^2$, les produits comme $(b_1b_2)^2$ ou $b_1^2b_2$ se réduisent à $b_1b_2$, puisque $b^2=b$ : avec deux bits seulement, tout polynôme est au plus de degré 2. Avec 3 bits par variable, $z_1=\sum_jc_jb_{1j}+c_0$ et $z_1z_2$ contient les produits $b_{1j}b_{2k}$ ; son carré contient $b_{1j}b_{1j'}b_{2k}b_{2k'}$ avec $j\ne j'$ et $k\ne k'$, de degré 4, que rien ne réduit.

\section*{Chapitre~\ref{ch:lagrangien} — Lagrangien}
\textbf{Ex.~\ref{ex:lag-saut}.} (a) Faisables : $A$ ($f=0$) et $B$ ($f=-0{,}2$) : l'optimum contraint est $B$. (b) $L_A=-2\lambda$, $L_B=-0{,}2-\lambda$, $L_C=-3+\lambda$, $L_D=5+3\lambda$. $L_B\le L_A$ exige $\lambda\le0{,}2$ ; $L_B\le L_C$ exige $\lambda\ge1{,}4$. C'est incompatible : $B$ n'est jamais le minimum. (c) $\lambda=0$ : $L=(0;\,-0{,}2;\,-3;\,5)$, les deux meilleures sont $C$ et $B$. $\lambda=2$ : $L=(-4;\,-2{,}2;\,-1;\,11)$, les deux meilleures sont $A$ et $B$. La réunion $\{A,B,C\}$ contient $B$, que la vérification retient (faisable, meilleur objectif). (d) Minimiser $f-\lambda h$ revient à pousser vers le haut une droite de pente $\lambda$ jusqu'à toucher un point : seuls les points de l'enveloppe convexe inférieure peuvent être touchés. $B=(1;\,-0{,}2)$ est au-dessus du segment $AC$ (qui vaut $-1$ en $h=1$) : il est invisible pour tout $\lambda$, mais il fait partie des bonnes lectures.

\textbf{Ex.~\ref{ex:lag-geo}.} (a) $B_i(0{,}5)=\binom4i/16$ : $\tfrac1{16},\tfrac4{16},\tfrac6{16},\tfrac4{16},\tfrac1{16}$. (b) $\Delta_0+4\Delta_1+6\Delta_2+4\Delta_3+\Delta_4\ge0$ avec $\Delta_i=A_{{\rm ext},i}-A_{{\rm int},i}$, chaque $\Delta_i$ étant affine en $z_i$ et $z_{i+5}$. (c) Les $B_i$ et $C$ sont positifs sur $[0,1]$ : une somme de termes positifs est positive. On ne l'a pas imposé parce que c'est une condition suffisante trop sévère : le lot L6 a montré que 3 des 10 profils du front de référence ont un $\Delta_i$ négatif tout en étant valides.

\section*{Chapitre~\ref{ch:classifieur} — Classifieur}
\textbf{Ex.~\ref{ex:clf-budget}.} (a) 93 échecs et 207 succès ; on évite $0{,}43\times93\approx40$ échecs et l'on sacrifie $0{,}063\times207\approx13$ succès. (b) 53 propositions rejetées ; remplacées par des propositions qui passent le classifieur, dont le taux de succès est $(207-13)/(300-53)\approx0{,}79$, elles donnent environ 42 succès : $194+42\approx236$ faisables au lieu de 207, soit 14~\% de plus. (c) Les succès sacrifiés ne se valent pas : si ce sont des profils du front (trois ont une probabilité inférieure à 0,2), on perd de l'hypervolume que ce bilan ne compte pas. C'est ce qu'on observe sur le problème V2.

\textbf{Ex.~\ref{ex:clf-seuil}.} $\ln(0{,}1/0{,}9)=-2{,}20$ ; $\ln(0{,}2/0{,}8)=-1{,}39$ ; $\ln1=0$. Un log-odds de $-1{,}8$ n'est accepté qu'au seuil 0,1. Sa probabilité de succès est $1/(1+e^{1{,}8})\approx0{,}14$.

\section*{Chapitre~\ref{ch:hors-ligne} — Sans XFOIL}
\textbf{Ex.~\ref{ex:hl-exact}.} Le QUBO a $2^{11}=2\,048$ chaînes (10 bits de conception, 1 auxiliaire) et chaque résolution exacte en garde les 600 meilleures, soit près de 30~\%. L'optimum contraint, qui a un bon objectif, figure dans cette liste pour tous les $\lambda$ visités, et la vérification le retient puisqu'elle trie selon l'objectif seul. À 30 bits, il faudrait garder une fraction comparable de $10^9$ chaînes : c'est l'énumération, impossible. Le recuit ne rend que quelques dizaines de lectures distinctes, et ce sont les $\lambda$ qui décident où elles tombent.

\textbf{Ex.~\ref{ex:hl-gain}.} (a) La boucle ne simule que la conception ; dans le QUBO à 14 qubits, la même conception occupe 16 chaînes, dont une seule de basse énergie. Comparer des chaînes complètes avec des références $1/2^{14}$ et $1/2^{10}$ fausserait la comparaison. (b) Un gain de 1 veut dire que QAOA ne tire pas la conception optimale plus souvent qu'un tirage uniforme. Il peut faire mieux ailleurs (80~\% d'échantillons faisables, contre bien moins au hasard), mais il n'aide pas à trouver l'optimum. (c) Le cours IA et QUBO mesurait la probabilité de l'état fondamental du QUBO complet, chaîne entière, rapportée à $1/2^{14}$ ; cet état n'est pas toujours l'optimum contraint (petites violations rentables), et la normalisation diffère.

\section*{Chapitre~\ref{ch:cqm} — Interface}
\textbf{Ex.~\ref{ex:cqm-ecrire}.}
\begin{lstlisting}[language={}]
Minimize
 obj: b0 - 2 b1 + [ 6 b0 * b1 ] / 2
Subject To
 c1: b0 + b1 >= 1
 c2: - 0.5 b0 + [ b0 * b1 ] >= -0.2
Binary
 b0 b1
End
\end{lstlisting}
Énumération : $00$ viole c1 ; $10$ : objectif 1, c2 vaut $-0{,}5<-0{,}2$, infaisable ; $01$ : objectif $-2$, faisable ; $11$ : objectif $2$, faisable. L'optimum est $01$. Le crochet de l'objectif est divisé par 2 : on y écrit $6$ pour un couplage de 3.
